\subsection{可测子空间上的诱导测度}

设 X 是 T 的子集，\(\nu\) 是映射 \(\mu: K \mapsto \mu_{K}\) 在 X 的紧子集全体上的限制；显然 \(\nu\) 是 X 上的一个前测度。另一方面，设 \(x \in X\) ，V 是 \(x\) 在 \(\mathbf{T}\) 中的开邻域且 \(\mu^{\bullet}(\mathrm{V}) < +\infty\) ；则

\[
\nu^{\bullet}(\mathrm{X}\cap \mathrm{V}) = \sup_{\substack{\mathrm{K\ compact}\\ \mathrm{K}\subset \mathrm{X}\cap \mathrm{V}}}\mu^{\bullet}(\mathrm{K})\leqslant \mu^{\bullet}(\mathrm{V}) <   + \infty
\]

（由 §1, No.~2 的注 3），因此 \(\nu\) 是一个测度。

当 X 不是 \(\mu\) -可测的时，负担 \(\nu^{\bullet}\) 与 \((\mu^{\bullet})_{X}\) 未必相等，测度 \(\nu\) 也并不令人感兴趣。

定义 1. --- 设 X 是 T 的 \(\mu\) -可测子集。\(\mu: \mathrm{K} \mapsto \mu_{\mathrm{K}}\) 在 X 的紧子集全体上的限制称为 \(\mu\) 在子空间 X 上诱导的测度，记作 \(\mu_{\mathrm{X}}\) 或 \(\mu|_{\mathrm{X}}\) 。

命题 1. --- 设 X 是 T 的 \(\mu\) -可测子集。负担 \((\mu_{\mathrm{X}})^{\bullet}\) 等于负担 \((\mu^{\bullet})_{\mathrm{X}}\) ——由 \(\mu^{\bullet}\) 在 X 上诱导——(§1, No.~1)。换言之，有 \((\mu_{\mathrm{X}})^{\bullet}(g) = \mu^{\bullet}(g^{0})\) （对每个函数 \(g \in \mathcal{F}_{+}(\mathrm{X})\)）。

设 \(f \in \mathcal{F}_{+}(X)\) ，\(f^{0}\) 是 \(f\) 到 \(T\) 上经 0 的延拓。有 \((\mu^{\bullet})_{X}(f) = \mu^{\bullet}(f^{0}) = \sup_{L} \mu^{\bullet}(f^{0}\varphi_{L})\) ，其中 \(L\) 取遍 \(T\) 的紧子集全体 (§1, No.~2, Prop. 2)；类似地 \((\mu_{X})^{\bullet}(f) = \sup_{K} \mu_{K}^{\bullet}(f_{K}) = \sup_{K} \mu^{\bullet}(f^{0}\varphi_{K})\) ，其中 \(K\) 取遍 \(X\) 的紧子集全体。于是问题归结为证明 \(\mu^{\bullet}(f^{0}\varphi_{L}) = \sup_{K} \mu^{\bullet}(f^{0}\varphi_{K})\) 对每个紧子集 \(L\) ⊂ \(T\) 成立，其中 \(K\) 取遍 \(L \cap X\) 的紧子集全体。现设 \((K_{n})\) 是含于 \(L \cap X\) 的递增紧集序列，使 \((L \cap X) - \bigcup_{n} K_{n}\) 局部 \(\mu\) -可忽略 (§1, No.~8, Prop. 11)；由于 \(f^{0}\) 在 \(X\) 外为零，\(f^{0}\varphi_{L}\) 在 \(L \cap X\) 外为零，故局部几乎处处等于序列 \((f^{0}\varphi_{K_{n}})\) 的上包络。由此推出 \(\mu^{\bullet}(f^{0}\varphi_{L}) = \sup \mu^{\bullet}(f^{0}\varphi_{K_{n}})\) ，从而得到所要的结果。

注 3). --- 设 X 是 T 的 \(\mu\) -可测子集。把 §1, No.~8 的 Prop. 10 应用于 \(g = \varphi_{\mathbf{X}}\) ，存在 T 的一个破碎 \((\mathrm{K}_{\alpha})_{\alpha \in \mathbf{A}}\) ，使对每个 \(\alpha \in \mathbf{A}\) ，或者 \(\mathrm{K}_{\alpha} \subset \mathrm{X}\) ，或者 \(\mathrm{K}_{\alpha} \subset \mathbb{C}\mathrm{X}\) 。若按 §1, No.~8 的附论的程序修改 T 的拓扑，则空间 \(\mathrm{X}'\) ——由给 X 赋予 \(\mathrm{T}'\) 诱导的拓扑而得到——是局部紧的，并且我们知道如何对 \(\mu\) （相应地对 \(\mu_{\mathbf{X}}\) ）联系一个测度 \(\mu'\) （相应地 \(\nu\) ），它在 \(\mathrm{T}'\) 上（相应地在 \(\mathrm{X}'\) 上），它与 \(\mu\) （相应地 \(\mu_{\mathbf{X}}\) ）有相同的本质上积分：这蕴含 \(\mu'_{\mathbf{X}'} = \nu\) 。由于局部可忽略集、可测映射以及以 Banach 空间为值的本性可积函数对 \(\mu\) 与 \(\mu'\) 相同，且对 \(\mu_{\mathbf{X}}\) 与 \((\mu_{\mathbf{X}})' = \nu = \mu'_{\mathbf{X}'}\) 相同，关于诱导测度的积分理论就化归为 Ch. V, §7 在局部紧空间这一特殊情形中处理过的理论。结果的改写留给读者。
% semantic-fragments: legacy-input-seam
\relax{}
